\documentclass[10pt,a4paper]{amsart}
\usepackage[margin=19mm]{geometry}
\usepackage{amsmath,amssymb,mathtools}
\usepackage[scr=rsfs,cal=euler]{mathalfa}
\usepackage{xfrac}
\usepackage[unicode,hidelinks,bookmarks=false]{hyperref}
\newcommand{\ie}{i.e.\ }
\newcommand{\cf}{cf.\ }
\newcommand{\resp}{respectively\ }
\newcommand{\Subsection}{\subsection}
\providecommand{\nobreakdash}{\nobreak\hbox{-}}
\setlength{\emergencystretch}{2em}
\multlinegap0pt
\usepackage{dsfont}
\usepackage{amssymb}
\usepackage{xcolor}
\usepackage{bm}
\usepackage{tikz}
\usepackage{float}
\newtheorem{lemma}{Lemma}
\newtheorem{theorem}{Theorem}
\def\R{\mathbb R}
\newcommand{\tr}{{\rm Tr}}
\title{A quadratic Grassmann manifold optimization problem arising from quantum embedding methods}
\author{Thomas Ayral, Eric Canc\`es, Fabian M. Faulstich, Lin Lin,, Alicia Negre}
\date{Source: arXiv:2603.17080v1 (CC BY 4.0)}
\begin{document}
\maketitle

\noindent\emph{Source and licence.} A quadratic Grassmann manifold optimization problem arising from quantum embedding methods. Version arXiv:2603.17080v1 (CC BY 4.0), distributed by arXiv under the Creative Commons Attribution 4.0 International licence, http://creativecommons.org/licenses/by/4.0/, as recorded in the arXiv licence metadata for this exact version and shown on the abstract page https://arxiv.org/abs/2603.17080v1. Full licence notice: \texttt{licenses/arxiv-2603.17080-CC-BY-4.0.txt}.\par\smallskip

\noindent\textit{Editorial scope.} Counted unit: the article's theorem thm:convex\_pbm (source0.tex lines 379--408). The article's complete original proof of it is delivered with it (source0.tex lines 410--456, one \texttt{proof} environment of 47 lines). No mathematics is written, repaired or re-derived: the statement and the proof are the article's own lines, in the article's own notation, and no retained line is substituted or re-flowed.  Retained uncounted as the source-local context the proof needs: lines 87--90; lines 96--99; lines 355--361.  Nothing was omitted from any retained span, so the package carries no omission record.  No retained line contains a citation, so the wrapper carries no bibliography and no bibliography entry.  No retained line contains a figure environment, an image-inclusion command or drawing code, so no image is delivered and the package declares no asset dependency.  Editorial formatting only, in addition: an AMS \texttt{amsart} preamble that restates the article's own declarations and macros used by the retained lines, namely the environments \texttt{lemma}, \texttt{theorem} on separate counters, together with the standard packages those lines need.  The retained environments print in this document as Lemma 1, Theorem 1 in order of appearance, whereas the article prints the same items with the numbering of its own full document.  The complete native source of the article as distributed in the arXiv e-print for this exact version is bundled unchanged as \texttt{sources/196599/source0.tex}.
\begin{equation}
\label{eq:OptProblem}
\min_{P \in \mathcal M} J(P), 
\end{equation}

\begin{equation}
\label{eq:CostFunctionJ}
    J(P) = \text{Tr}(BP)- \frac{1}{2} \text{Tr}(A P A P) = \text{Tr}(BP) - \frac{1}{2} \lVert A^{\frac{1}{2}} P A^{\frac{1}{2}} \rVert^2 \quad \mbox{with} \quad A,B \in \R^{M \times M}_{\rm sym}, \quad A \succeq 0,
\end{equation}

\begin{lemma} \label{lem:convexification}
    Let $\widetilde J:\R^{M \times M}_{\rm sym} \to \R$ be the quadratic function defined by
\begin{equation}\label{eq:def_tildeJ}
\widetilde J(D) := \tr(CD) + \frac 14 \| [A,D] \|^2  \quad \mbox{with} \quad C := B- \frac 12 A^2.
\end{equation}
Then, for all $P \in \mathcal M$, $\widetilde J(P)=J(P)$. In particular, $J$ and $\widetilde J$ have the same critical points on $\mathcal M$.
\end{lemma}

\begin{theorem}
\label{thm:convex_pbm}
Let $A \in \R^{M \times M}_{\rm sym}$ be a positive semidefinite matrix, and $C \in \R^{M \times M}_{\rm sym}$.
The set $\mathcal D_\star$ of the (global) minimizers of the convex optimization problem
\begin{equation} \label{eq:convex_opt}
   \widetilde I= \min_{D \in \mathcal K} \widetilde J(D)
\end{equation}
is a non-empty, compact, convex subset of $\mathcal K$, and we have the following properties:
\begin{enumerate}
    \item The function $\R^{M\times M}_{\rm sym} \ni D \mapsto \nabla \widetilde J(D)=C- \frac 12[[A,D],A] \in \R^{M \times M}_{\rm sym}$ takes a constant value $H_\star$ on $\mathcal D_\star$.
    \item Let $\mu_1 \le \cdots \le \mu_M$ be the eigenvalues of $H_\star$, counted with their multiplicities and ranked in non-decreasing order. Moreover, define 
    \begin{equation*}
        \widetilde {\mathcal D}_\star:= \left\{ D \in \mathcal K \; | \;  \mathds 1_{(-\infty,\mu_m)}(H_\star) \preceq D \preceq 1_{(-\infty,\mu_m]}(H_\star) \right\}.
    \end{equation*}
    Then
    \begin{equation} \label{eq:DsubsetDtilde}
        {\mathcal D}_\star \subset \widetilde {\mathcal D}_\star,
    \end{equation}
    and 
    \begin{align}
    \forall D_\star \in \mathcal D_\star, \quad \forall \widetilde D_\star \in \widetilde{\mathcal D}_\star, \quad \widetilde D_\star \in \mathcal D_\star \Leftrightarrow \;  [A,D_\star-\widetilde D_\star]=0 \;   \Rightarrow  \;  {\rm Ran}(D_*-\widetilde D_*) \subset \mathcal V, \label{eq:charac_D_star}
    \end{align}
    where $\mathcal V \subset \R^M$ is the largest $A$-invariant vector subspace of ${\rm Ker}(H_\star-\mu_m)$.
\item If $\mu_m < \mu_{m+1}$, then 
\begin{equation}
    P_\star:= \mathds 1_{(-\infty,\mu_m]}(H_\star)
\end{equation}
is the unique minimizer of \eqref{eq:convex_opt}, i.e., $\mathcal D_\star=\{P_\star\}$, and $P_\star$ is also the unique global minimizer of the original problem~\eqref{eq:OptProblem}--\eqref{eq:CostFunctionJ}.
\end{enumerate}
\end{theorem}

\begin{proof}
    The function $\widetilde J$ is convex on $\R^{M \times M}_{\rm sym}$ and $\mathcal K$ is a compact convex subset of $\R^{M \times M}_{\rm sym}$. The set $\mathcal{D}_\star$  of minimizers of problem \eqref{eq:convex_opt} therefore is a non-empty compact convex subset of $\mathcal{K}$. Let $D_\star \in \mathcal{D}_\star$ be a (global) minimizer of~\eqref{eq:convex_opt}. Using convexity, one has, for any $D \in \mathcal{K}$,
\begin{equation}\label{eq:optimal_0}
\forall \theta \in [0,1], \quad \widetilde J((1-\theta) D_\star + \theta D) \geq \widetilde J(D_\star). \end{equation}
For small $\theta$, $\theta \geq 0$, one can expand \eqref{eq:optimal_0} to get
\begin{align*}
    \widetilde J((1-\theta) D_\star + \theta D)& = \widetilde J( D_\star + \theta (D-D_\star))= \widetilde J( D_\star) + \theta \langle\nabla \widetilde J( D_\star) , D-D_\star\rangle + o(\theta),
\end{align*}
yielding the Euler inequality
\begin{equation}\label{eq:euler_lagrange}
\forall D \in \mathcal{K}, \quad \tr(H_\star D) \ge  \tr( H_\star D_\star), \quad \mbox{with} \quad H_\star := \nabla \widetilde J( D_\star)=C- \frac 12[[A,D_\star],A].
\end{equation}
Denoting by  $\mu_1 \le \cdots \le \mu_M$ the eigenvalues of $H_\star$, we have 
$$
    \min_{D \in \mathcal{K} }\tr(H_\star D) = \min_{D \in \mathcal{M}}\tr(H_\star D) = \sum_{k = 1}^{m} \mu_{k},
$$
and
$$
D_\star \in \widetilde {\mathcal D}_\star:= \mathop{\rm argmin}_{D \in \mathcal K}\tr(H_\star D)=\left\{ D \in \mathcal K \; | \;  \mathds 1_{(-\infty,\mu_m)}(H_\star) \preceq D \preceq 1_{(-\infty,\mu_m]}(H_\star) \right\}.
$$
Let $\widetilde D_\star$ be another minimizer of~\eqref{eq:convex_opt}. By convexity, we have that
\begin{align}
\forall \theta \in [0,1], \quad \widetilde I &= \widetilde J((1-\theta) D_\star + \theta \widetilde D_\star)  = \widetilde J(D_\star +\theta(\widetilde D_\star- D_\star)) \nonumber \\
&= \widetilde I + \theta \tr(H_\star ( \widetilde D_\star-D_\star)) +\frac{\theta^2}4 \lVert [A,\widetilde D_{\star} - D_\star] \rVert^2. \label{eq:sec_order_conv}
\end{align}
This implies $\tr(H_\star ( \widetilde D_\star-D_\star))=0$ which yields $\widetilde D_* \in \widetilde {\mathcal D}_\star$ and therewith \eqref{eq:DsubsetDtilde}. Moreover, \eqref{eq:sec_order_conv} implies $[A,\widetilde D_{\star} - D_\star]=0$, which yields
$$
\nabla \widetilde J(\widetilde D_*) = C-  \frac 12[[A,\widetilde D_\star],A] =  C-  \frac 12[[A,D_\star],A] = H_\star,
$$
proving the first assertion of the theorem.

\medskip

Let $D_\star \in \mathcal D_\star$ and $\widetilde D_\star \in \widetilde{\mathcal D}_\star$. As $\tr(H_\star D_\star) = \tr(H_\star \widetilde D_\star)$, we deduce from~\eqref{eq:sec_order_conv} that $\widetilde D_\star \in \mathcal D_\star$ if and only if $[A,D_\star-\widetilde D_\star]=0$. This proves the equivalence in~\eqref{eq:charac_D_star}. To prove the implication, we observe that $\widetilde D_\star$ is in $\widetilde{\mathcal D}_\star$ if and only if ${\rm Ran}(\widetilde D_\star-D_\star) \subset {\rm Ker}(H_\star-\mu_m)$ and $[A,\widetilde D_\star-D_\star]=0$, that is if and only if there exists a family $(u_1,\cdots,u_d)$ of vectors in ${\rm Ker}(H_\star-\mu_m)$ such that
$$
\widetilde D_\star-D_\star = \sum_{j=1}^{d} n_j u_ju_j^T 
$$
with
$$
d \le {\rm dim}({\rm Ker}(H_\star-\mu_m)), \quad n_j \in [-1,1] \setminus \{0\}, \quad \sum_{j=1}^d n_j=0, \quad Au_j= \alpha_j u_j, \quad u_j^Tu_k=\delta_{jk}.
$$
This condition implies that ${\rm Ran}(D_*-\widetilde D_*) = \mbox{Span}(u_1,\cdots,u_d) \subset \mathcal V$, where $\mathcal V \subset \R^M$ is the largest $A$-invariant vector subspace of ${\rm Ker}(H_\star-\mu_m)$. This completes the proof of the second assertion.

\medskip

The third assertion readily follows from~\eqref{eq:DsubsetDtilde}, the fact that if $\mu_m < \mu_{m+1}$, then $\widetilde {\mathcal D}_\star$ is the singleton $\{1_{(-\infty,\mu_m](H_\star)}\}$, and Lemma~\ref{lem:convexification}.
\end{proof}

\end{document}
